\documentclass[../../main.tex]{subfiles}

\begin{document}

\chapter{Operasi Logika dan Shift}
\label{chap:logika}

% =========================================================
% OUTCOME PEMBELAJARAN KHUSUS BAB
% =========================================================

\begin{subcpmkbox}
\vspace{1.1em}
\SubCPMK{3.2} --- Menerapkan operasi NOT, AND, OR, XOR, TEST, SHL,
dan SHR untuk manipulasi bit.

\vspace{0.4em}
\textbf{CPMK yang didukung:} \CPMK{03} \quad\textbullet\quad
\textbf{CPL yang didukung:} \CPL{03}
\end{subcpmkbox}

\begin{tujuanbox}
\vspace{1.1em}
Setelah mempelajari bab ini, mahasiswa mampu:
\begin{enumerate}
    \item menjelaskan tabel kebenaran \asm{AND}, \asm{OR},
          \asm{XOR}, dan \asm{NOT};
    \item menjelaskan pengaruh operasi logika terhadap
          \eng{flag};
    \item menyalakan, memadamkan, membalik, dan menguji bit
          tertentu dengan \eng{mask};
    \item membedakan \asm{TEST} dan \asm{CMP};
    \item menggunakan \asm{SHL} dan \asm{SHR} untuk menggeser bit;
    \item memakai pergeseran untuk perkalian dan pembagian dengan
          pangkat dua;
    \item menyusun program manipulasi bit seperti pemeriksaan
          paritas dan \eng{XOR cipher}.
\end{enumerate}
\end{tujuanbox}

\begin{materibox}
\vspace{1.1em}
\begin{enumerate}
    \item Operasi \asm{NOT}, \asm{AND}, \asm{OR}, dan \asm{XOR}.
    \item Pengaruh operasi logika terhadap \eng{flag}.
    \item \eng{Mask} dan manipulasi bit tertentu.
    \item \asm{TEST} sebagai penguji bit.
    \item \asm{SHL} dan \asm{SHR}.
    \item Perkalian dan pembagian dengan pangkat dua.
    \item Penerapan: paritas dan \eng{XOR cipher}.
\end{enumerate}
\end{materibox}

\subfile{section-01}
\subfile{section-02}
\subfile{section-03}

\section{Contoh}
\label{sec:contoh-bab11}

Dua contoh berikut memperlihatkan manipulasi bit pada program utuh.
Contoh pertama memakai \asm{SHR} dan \asm{TEST} untuk memeriksa
paritas. Contoh kedua memakai \asm{XOR} untuk enkripsi dan dekripsi
sederhana.

\subsection{Contoh 1: Memeriksa Paritas}

Program pada Kode~\ref{lst:paritas} menghitung jumlah bit 1 pada
sebuah byte, lalu menentukan paritasnya.

\lstinputlisting[
    language=MASM,
    caption={Program \file{paritas.asm} memeriksa paritas byte.},
    label={lst:paritas}
]{\CodeDir/paritas.asm}

Bagian penting program:

\begin{itemize}
    \item \lstasm|MOV CX, 8| menyiapkan delapan kali pergeseran.
    \item \lstasm|SHR AL, 1| menggeser bit terkanan ke \flag{CF}.
    \item \lstasm|ADC BL, 0| menambahkan \flag{CF} ke penghitung,
          sehingga \reg{BL} berisi jumlah bit 1.
    \item \lstasm|TEST BL, 1| memeriksa bit terendah jumlah bit.
          Bila bit itu 1, jumlahnya ganjil; bila 0, jumlahnya genap.
\end{itemize}

Byte uji \bin{10110110} memiliki lima bit 1, sehingga paritasnya
ganjil. Keluaran program tampak pada
Kode~\ref{lst:keluaran-paritas}.

\begin{lstlisting}[
    style=sessiongaya,
    caption={Keluaran program \file{paritas.asm}.},
    label={lst:keluaran-paritas}
]
Byte 10110110b berparitas: GANJIL
\end{lstlisting}

\subsection{Contoh 2: XOR Cipher}

Program pada Kode~\ref{lst:xor-cipher} memakai \asm{XOR} untuk
mengenkripsi dan mendekripsi string. Sifat $A \oplus K \oplus K = A$
membuat proses dapat dibalik dengan kunci yang sama.

\lstinputlisting[
    language=MASM,
    caption={Program \file{xor-cipher.asm} mengenkripsi dan
    mendekripsi string dengan XOR.},
    label={lst:xor-cipher}
]{\CodeDir/xor-cipher.asm}

Bagian penting program:

\begin{itemize}
    \item \lstasm|MOV AL, [SI]| mengambil satu byte dari string asli.
    \item \lstasm|XOR AL, kunci| membalik bit sesuai kunci.
    \item \lstasm|MOV [DI], AL| menyimpan hasil ke penyangga.
    \item Pemeriksaan \lstasm|CMP AL, '$'| menghentikan perulangan
          pada akhir string.
    \item Dekripsi memakai \asm{XOR} yang sama pada hasil enkripsi.
\end{itemize}

Dengan kunci \hex{20}, huruf besar berubah menjadi huruf kecil
karena bit ke-5 dibalik. Keluaran program tampak pada
Kode~\ref{lst:keluaran-xor}.

\begin{lstlisting}[
    style=sessiongaya,
    caption={Keluaran program \file{xor-cipher.asm}.},
    label={lst:keluaran-xor}
]
Teks asli    : BELAJARASM
Terenkripsi  : belajarasm
Terdekripsi  : BELAJARASM
\end{lstlisting}

\begin{infobox}
\vspace{1.1em}
\eng{XOR cipher} pada contoh ini hanya demonstrasi sifat \asm{XOR},
bukan algoritma kriptografi yang aman. Kunci satu byte dan pola
sederhana mudah ditebak; pembahasan keamanan berada di luar cakupan
buku ini.
\end{infobox}

\section{Praktikum}
\label{sec:praktikum-bab11}

\begin{praktikumbox}
\vspace{1.1em}
\textbf{Judul:} Manipulasi bit dengan SHR, TEST, dan XOR.\\[0.3em]
\textbf{Sub-CPMK:} 3.2 \quad\textbullet\quad
\textbf{Alokasi waktu:} 100 menit \quad\textbullet\quad
\textbf{Bentuk:} individu
\end{praktikumbox}

\subsection{Tujuan Praktikum}

Setelah menyelesaikan praktikum ini, mahasiswa mampu menggunakan
\asm{SHR}, \asm{TEST}, dan \asm{XOR} untuk memeriksa serta mengubah
bit, lalu menjelaskan hasilnya berdasarkan isi register dan
\eng{flag}.

\subsection{Alat dan Bahan}

\begin{itemize}
    \item Komputer dengan emu8086 terpasang.
    \item DOSBox dengan TASM, TLINK, dan DEBUG.
    \item Berkas \file{paritas.asm} dan \file{xor-cipher.asm} dari
          \file{code/asm/}.
    \item Lembar pengamatan register dan \eng{flag}.
\end{itemize}

\subsection{Langkah Kerja}

\begin{enumerate}
    \item \textbf{Menguji program paritas.}
    \begin{enumerate}
        \item Rakit \file{paritas.asm} dengan TASM, lalu link
              dengan TLINK.
        \item Jalankan program dan catat keluarannya.
        \item Buka program pada emu8086, lalu jalankan
              \eng{single step} sampai perulangan selesai.
        \item Catat \reg{AL}, \reg{BL}, \reg{CX}, dan \flag{CF}
              pada setiap pergeseran.
        \item Jelaskan mengapa \reg{BL} berisi jumlah bit 1.
    \end{enumerate}

    \item \textbf{Mengubah paritas.}
    \begin{enumerate}
        \item Ubah \code{nilai} menjadi \bin{10110111} yang memiliki
              enam bit 1.
        \item Rakit ulang dan jalankan program.
        \item Pastikan keluaran berubah menjadi \code{GENAP}.
    \end{enumerate}

    \item \textbf{Menguji XOR cipher.}
    \begin{enumerate}
        \item Rakit dan jalankan \file{xor-cipher.asm}.
        \item Catat teks asli, hasil enkripsi, dan hasil dekripsi.
        \item Jelaskan mengapa hasil dekripsi sama dengan teks asli.
    \end{enumerate}

    \item \textbf{Mengubah kunci XOR.}
    \begin{enumerate}
        \item Ubah \code{kunci} dari \hex{20} menjadi \hex{01}.
        \item Rakit ulang, lalu jalankan program.
        \item Catat hasil enkripsi dan dekripsi.
        \item Jelaskan mengapa hasil enkripsi tidak lagi terbaca
              sebagai huruf.
    \end{enumerate}

    \item \textbf{Menguji mask.}
    \begin{enumerate}
        \item Tulis potongan program yang menyalakan bit 3 pada
              \reg{AL} = \hex{A5}, lalu memadamkannya kembali.
        \item Gunakan \asm{OR} dan \asm{AND} dengan mask yang tepat.
        \item Periksa hasilnya dengan \eng{single step}.
    \end{enumerate}

    \item \textbf{Memeriksa dengan DEBUG.}
    \begin{enumerate}
        \item Jalankan \file{debug paritas.exe}.
        \item Ketik \file{r} untuk melihat register.
        \item Ketik \file{t} beberapa kali dan catat perubahan
              \reg{AL}, \reg{BL}, dan \flag{CF}.
    \end{enumerate}
\end{enumerate}

\subsection{Hasil yang Dikumpulkan}

\begin{enumerate}
    \item Berkas \file{paritas.asm} dan \file{xor-cipher.asm} versi
          modifikasi.
    \item Tangkapan layar keluaran kedua program.
    \item Tabel pengamatan \reg{AL}, \reg{BL}, \reg{CX}, dan
          \flag{CF} pada perulangan paritas.
    \item Penjelasan sifat $A \oplus K \oplus K = A$.
    \item Catatan hasil percobaan mask.
\end{enumerate}

\begin{asesmenbox}
\vspace{1.1em}
\textbf{Asesmen praktikum 9 (bobot 5\% dari komponen praktikum).}

\begin{tabular}{@{}p{6.2cm}cp{4.6cm}@{}}
    \toprule
    \textbf{Indikator} & \textbf{Bobot} & \textbf{Kriteria} \\
    \midrule
    Ketepatan pengamatan pergeseran & 30\% &
    \reg{AL}, \reg{BL}, \reg{CF} tercatat benar \\
    \addlinespace
    Ketepatan analisis paritas & 25\% &
    Hubungan jumlah bit 1 dan \asm{TEST} dijelaskan \\
    \addlinespace
    Ketepatan analisis XOR & 25\% &
    Enkripsi dan dekripsi dijelaskan dengan sifat XOR \\
    \addlinespace
    Ketepatan penggunaan mask & 20\% &
    Mask dan instruksi yang dipilih benar \\
    \bottomrule
\end{tabular}
\end{asesmenbox}

\section{Latihan}
\label{sec:latihan-bab11}

\begin{enumerate}
    \item \textbf{Tabel kebenaran.} Jelaskan hasil \asm{AND},
          \asm{OR}, \asm{XOR}, dan \asm{NOT} untuk pasangan bit
          berikut.
          \begin{enumerate}
              \item $A = 1$, $B = 0$
              \item $A = 1$, $B = 1$
              \item $A = 0$, $B = 0$
          \end{enumerate}

    \item \textbf{Flag.} Jelaskan pengaruh \asm{AND}, \asm{OR}, dan
          \asm{XOR} terhadap \flag{CF}, \flag{OF}, \flag{ZF}, dan
          \flag{SF}. Mengapa \asm{NOT} berbeda?

    \item \textbf{Mask.} Diketahui \reg{AL} = \hex{A5}. Tuliskan
          instruksi untuk:
          \begin{enumerate}
              \item menyalakan bit 4;
              \item memadamkan bit 4;
              \item membalik bit 4;
              \item menguji bit 4 tanpa mengubah \reg{AL}.
          \end{enumerate}

    \item \textbf{TEST dan CMP.} Jelaskan perbedaan \asm{TEST} dan
          \asm{CMP}. Kapan masing-masing sebaiknya dipakai?

    \item \textbf{TEST AL, AL.} Apa fungsi \lstasm|TEST AL, AL|?
          Mengapa instruksi ini dipakai untuk memeriksa apakah
          \reg{AL} bernilai nol?

    \item \textbf{Pergeseran.} Diketahui \reg{AL} = \bin{10100101}.
          Tentukan hasil dan \flag{CF} untuk:
          \begin{enumerate}
              \item \lstasm|SHL AL, 1|
              \item \lstasm|SHR AL, 1|
              \item \lstasm|SHL AL, 2|
          \end{enumerate}

    \item \textbf{Perkalian dan pembagian.} Tuliskan instruksi untuk
          menghitung $5 \times 16$ dan $1000 \div 8$ dengan
          pergeseran. Jelaskan mengapa cara ini hanya tepat untuk
          bilangan tidak bertanda.

    \item \textbf{Analisis program.} Pada \file{paritas.asm},
          jelaskan mengapa \lstasm|ADC BL, 0| dipakai setelah
          \lstasm|SHR AL, 1|. Pada \file{xor-cipher.asm}, jelaskan
          mengapa \asm{XOR} yang sama dapat dipakai untuk enkripsi
          dan dekripsi.
\end{enumerate}

\begin{refleksibox}
\vspace{1.1em}
Jawab pertanyaan berikut pada jurnal praktikum Anda.
\begin{enumerate}
    \item Operasi mana yang paling sulit Anda pahami: mask,
          \asm{TEST}, atau pergeseran? Mengapa?
    \item Setelah memeriksa \flag{CF} pada pergeseran, apakah Anda
          dapat menjelaskan ke mana bit yang keluar? Jelaskan.
    \item Konsep apa yang perlu Anda ulangi sebelum mempelajari
          mode pengalamatan pada bab berikutnya?
\end{enumerate}
\end{refleksibox}

\section{Rangkuman}
\label{sec:rangkuman-bab11}

Pada bab ini telah dibahas:

\begin{itemize}
    \item Operasi logika bekerja pada setiap bit secara mandiri,
          tanpa menjalarkan \eng{carry} antarbit.

    \item \asm{AND} menghasilkan 1 bila kedua bit bernilai 1,
          \asm{OR} bila salah satu bit bernilai 1, \asm{XOR} bila
          kedua bit berbeda, dan \asm{NOT} membalik seluruh bit.

    \item \asm{AND}, \asm{OR}, dan \asm{XOR} selalu mengosongkan
          \flag{CF} dan \flag{OF}. \flag{ZF}, \flag{SF}, dan
          \flag{PF} mengikuti hasil. \asm{NOT} tidak mengubah
          \eng{flag}.

    \item \eng{Mask} dipakai untuk menyalakan bit dengan \asm{OR},
          memadamkan bit dengan \asm{AND} dan mask terbalik,
          membalik bit dengan \asm{XOR}, serta menguji bit dengan
          \asm{TEST}.

    \item \asm{TEST} melakukan \asm{AND} tanpa menyimpan hasil,
          sedangkan \asm{CMP} melakukan \asm{SUB} tanpa menyimpan
          hasil. \asm{TEST} dipakai untuk bit, \asm{CMP} untuk nilai
          keseluruhan.

    \item \asm{SHL} menggeser bit ke kiri dan mengisi bit kanan
          dengan 0. \asm{SHR} menggeser bit ke kanan dan mengisi bit
          kiri dengan 0. Bit yang keluar ditampung \flag{CF}.

    \item Pergeseran ke kiri setara dengan perkalian pangkat dua,
          sedangkan pergeseran ke kanan setara dengan pembagian
          pangkat dua untuk bilangan tidak bertanda.

    \item Penerapan operasi logika dan pergeseran meliputi
          pemeriksaan paritas, manipulasi bit, dan \eng{XOR cipher}
          sederhana.
\end{itemize}

\begin{checklistbox}
\vspace{1.1em}
Tandai setiap butir yang sudah Anda kuasai.

\begin{itemize}
    \item[\checkbox] Saya dapat menjelaskan tabel kebenaran
          \asm{AND}, \asm{OR}, \asm{XOR}, dan \asm{NOT}.
    \item[\checkbox] Saya dapat menjelaskan pengaruh operasi logika
          terhadap \eng{flag}.
    \item[\checkbox] Saya dapat menyalakan, memadamkan, membalik,
          dan menguji bit dengan mask.
    \item[\checkbox] Saya dapat membedakan \asm{TEST} dan
          \asm{CMP}.
    \item[\checkbox] Saya dapat menggunakan \asm{SHL} dan
          \asm{SHR}.
    \item[\checkbox] Saya dapat memakai pergeseran untuk perkalian
          dan pembagian pangkat dua.
    \item[\checkbox] Saya dapat menjelaskan program paritas dan
          \eng{XOR cipher}.
\end{itemize}

\vspace{0.4em}
Bila masih ada butir yang belum bertanda, ulangi bagian yang
bersesuaian sebelum melanjutkan ke bab berikutnya.
\end{checklistbox}

\begin{infobox}
\vspace{1.1em}
\textbf{Rujukan bab.} Materi pada bab ini diringkas dari
\cite{gunawan2019}, Bagian 7, hlm.~68--73, dan \cite{sto2001},
BAB X, hlm.~47--51.
\end{infobox}

\section{Evaluasi}
\label{sec:evaluasi-bab11}

\subsection{Soal Pilihan Ganda}

\begin{enumerate}
    \item Nilai \flag{CF} dan \flag{OF} setelah instruksi
          \lstasm|XOR AL, BL| dijalankan adalah \dots
          \begin{enumerate}
              \item keduanya 0
              \item keduanya 1
              \item \flag{CF} = 1 dan \flag{OF} = 0
              \item bergantung pada nilai \reg{AL} dan \reg{BL}
          \end{enumerate}

    \item Instruksi untuk memadamkan bit ke-3 \reg{AL} tanpa mengubah
          bit lain adalah \dots
          \begin{enumerate}
              \item \lstasm|OR AL, 08h|
              \item \lstasm|AND AL, 08h|
              \item \lstasm|AND AL, 0F7h|
              \item \lstasm|XOR AL, 08h|
          \end{enumerate}

    \item Diketahui \reg{AL} = \bin{1111 1111} yang ditafsirkan
          sebagai $-1$. Instruksi \lstasm|SHR AL, 1| menghasilkan
          nilai \dots
          \begin{enumerate}
              \item $-1$ karena bit tanda dipertahankan
              \item $+127$ karena bit tanda diisi nol
              \item $0$ karena seluruh bit bergeser keluar
              \item $-64$ karena hasilnya dibulatkan ke bawah
          \end{enumerate}

    \item Instruksi yang menyertakan \flag{CF} sebagai bit tambahan
          yang ikut berputar adalah \dots
          \begin{enumerate}
              \item \asm{ROL} dan \asm{ROR}
              \item \asm{RCL} dan \asm{RCR}
              \item \asm{SHL} dan \asm{SHR}
              \item \asm{SAL} dan \asm{SAR}
          \end{enumerate}

    \item Perbedaan hakiki antara \asm{TEST} dan \asm{CMP} terletak
          pada \dots
          \begin{enumerate}
              \item operasi internalnya, yaitu \asm{AND} dan \asm{SUB},
                    meskipun keduanya tidak menyimpan hasil
              \item \asm{TEST} menyimpan hasil sedangkan \asm{CMP}
                    tidak
              \item \asm{TEST} hanya dapat dipakai pada operand 8 bit
              \item keduanya mengubah \flag{CF}, tetapi hanya
                    \asm{CMP} yang mengubah \flag{ZF}
          \end{enumerate}
\end{enumerate}

\subsection{Soal Uraian}

\begin{enumerate}
    \item Jelaskan mengapa \flag{CF} dan \flag{OF} selalu 0 setelah
          operasi logika, sedangkan pada operasi aritmatika keduanya
          dapat bernilai 1. Kaitkan dengan penjalaran \eng{carry}
          antarbit.

    \item Diketahui \reg{AL} = \bin{1100 1010}. Tuliskan hasil
          \lstasm|SHL AL, 1|, \lstasm|SHR AL, 1|, dan
          \lstasm|SAR AL, 1| beserta nilai \flag{CF}, lalu nyatakan
          hasil \asm{SAR} sebagai bilangan bertanda.

    \item Tuliskan program yang menyusun digit puluhan dari variabel
          \code{puluhan} dan digit satuan dari variabel \code{satuan}
          menjadi satu byte BCD, lalu memisahkannya kembali. Jelaskan
          peran \asm{SHL}, \asm{OR}, \asm{AND}, dan \asm{SHR}.

    \item Jelaskan mengapa dekripsi \eng{XOR cipher} dapat memakai
          kunci dan operasi yang sama dengan enkripsi, lalu sebutkan
          dua alasan mengapa cara ini belum memadai sebagai pengamanan
          data.
\end{enumerate}

\subsection{Kisi-kisi Penilaian}

\begin{table}[H]
    \centering
    \caption{Kisi-kisi evaluasi Bab~\ref{chap:logika}.}
    \label{tab:kisi-bab11}
    \small
    \setlength{\tabcolsep}{5pt}
    \begin{tabular}{@{}p{1.5cm}p{5.4cm}p{3.6cm}c@{}}
        \toprule
        \textbf{Sub-CPMK} & \textbf{Indikator yang diukur} &
        \textbf{Bentuk soal} & \textbf{Bobot} \\
        \midrule
        3.2 & Menjelaskan tabel kebenaran dan pengaruh operasi logika
        terhadap \eng{flag} & Pilihan ganda 1; Uraian 1 & 20\% \\
        \addlinespace
        3.2 & Menerapkan manipulasi bit ke-$n$ dan memilih instruksi
        yang tepat & Pilihan ganda 2, 5 & 20\% \\
        \addlinespace
        3.2 & Menganalisis pergeseran logis dan aritmetika pada
        bilangan bertanda & Pilihan ganda 3; Uraian 2 & 25\% \\
        \addlinespace
        3.2 & Menjelaskan rotasi dan peran \flag{CF} & Pilihan ganda
        4 & 10\% \\
        \addlinespace
        3.2 & Menulis program \eng{packing} BCD dan menganalisis
        \eng{XOR cipher} & Uraian 3, 4 & 25\% \\
        \bottomrule
    \end{tabular}
\end{table}

\begin{checklistbox}
\textbf{Checklist Pencapaian Kompetensi Bab~\ref{chap:logika}}

Seorang mahasiswa dinyatakan kompeten pada Sub-CPMK 3.2 apabila:
\begin{itemize}
    \item[\checked] dapat menyusun tabel kebenaran \asm{AND},
          \asm{OR}, \asm{XOR}, dan \asm{NOT};
    \item[\checked] dapat menjelaskan pengaruh operasi logika terhadap
          \flag{CF}, \flag{OF}, \flag{ZF}, \flag{SF}, dan \flag{PF};
    \item[\checked] dapat menuliskan instruksi penyalakan, pemadaman,
          pembalikan, dan pengujian bit ke-$n$;
    \item[\checked] dapat membedakan \asm{TEST} dan \asm{CMP} beserta
          pemakaiannya;
    \item[\checked] dapat menentukan hasil \asm{SHL}, \asm{SHR}, dan
          \asm{SAR} pada bilangan bertanda disertai \flag{CF};
    \item[\checked] dapat menjelaskan perbedaan \asm{ROL}/\asm{ROR}
          dengan \asm{RCL}/\asm{RCR};
    \item[\checked] berhasil menulis dan menjalankan program paritas
          serta \eng{packing} BCD.
\end{itemize}
\end{checklistbox}


\end{document}
